\documentclass[12pt]{article}                                                                                                           
\usepackage[utf8]{inputenc}                                                                                                             
\usepackage[T1]{fontenc}                                                                                                                
\usepackage{amsmath, amssymb}                                                                                                           
\usepackage{geometry}                                                                                                                   
\usepackage{xcolor}                                                                                                                     
\usepackage{lipsum}                                                                                                                     
\usepackage{titlesec}                                                                                                                   
\usepackage{fancyhdr}                                                                                                                   
\usepackage[most]{tcolorbox}                                                                                                                  
\usepackage{graphicx}                                                                                                                   
\usepackage{float}                                                                                                      
\usepackage{tikz}
\usetikzlibrary{shapes.geometric, arrows.meta}
\usetikzlibrary{arrows.meta, decorations.markings}                                                                                      
                                                                                                                                        
\geometry{a4paper, margin=2.5cm}                                                                                                        
\pagestyle{fancy}                                                                                                                       
\fancyhf{}                                                                                                                              
\rhead{Problem Set IV}                                                                                                                   
\lhead{Rigid body mechanics}                                                                                                                     
\cfoot{\thepage}                                                                                                                        
                                                                                                                                        
% Fancy titles                                                                                                                          
\titleformat{\section}{\normalfont\Large\bfseries}{Exercise \thesection:}{1em}{}                                                        
\titleformat{\subsection}{\normalfont\bfseries}{Correction:}{1em}{}                                                                     
                                                                                                                                        
% Custom box for correction with breakable option
\newtcolorbox{correctionbox}{                                                                                                           
  colback=gray!5,                                                                                                                       
  colframe=black,                                                                                                                       
  fonttitle=\bfseries,                                                                                                                  
  title=Correction,                                                                                                                     
  before skip=10pt,                                                                                                                     
  after skip=10pt,
  breakable,
  enhanced
}                                                                                                                                       
                                                                                                                                        
% Fix headheight warning
\setlength{\headheight}{14.49998pt}

\begin{document}

\begin{center}
\Large\textbf{Ibn Tofail University} \\[1em]
\large\textit{Rigid body mechanics} \\[2em]
\large\textit{Problem Set IV} \\[0.5em]
\end{center}

\vspace{1cm}
% ----------- EXERCISE 1 -------------
\section{}

A perfect rigid body having the shape of a square $ABCD$ is moved from position $A_1B_1C_1D_1$ to position $A_2B_2C_2D_2$.

\begin{enumerate}
    \item Show, using a figure, that the motion of this solid can be considered as the sum of a translation and a rotation about a suitably chosen fixed point.
    \item Give the value of the angle of rotation.
\end{enumerate}

\begin{correctionbox}

\end{correctionbox}

% ----------- EXERCISE 2 -------------
\section{}

Let $R(O, X, Y, Z)$ be a reference frame assumed fixed. A homogeneous rod of length $L$ and mass $m$ rotates, in the vertical plane $(XOY)$, about the axis $OZ$. The position of the rod with respect to $R$ is indicated by the angle $\theta$ (See Figure below).

Determine the kinetic elements of the solid, at point $O$, with respect to frame $R$.

\begin{correctionbox}

\end{correctionbox}

% ----------- EXERCISE 3 -------------
\section{}

A homogeneous and non-deformable hoop of radius $R$, center $G$, and mass $m$, rolls without slipping on an inclined plane of angle $\alpha$.

We denote by $x$ the abscissa of $G$ and by $\theta$ the angle that indicates the position of a point on the hoop. Initially, the geometric contact point coincides with point $O$ and $\theta = 0$.

\begin{enumerate}
    \item What is the number of degrees of freedom of the hoop?
    \item Determine the kinetic elements of the hoop at point $O$.
\end{enumerate}

\begin{correctionbox}

\end{correctionbox}

% ----------- EXERCISE 4 -------------
\section{}

A homogeneous solid cone $S$, of mass $m$, radius $R$, and height $h$ is in motion about its fixed apex $O$, relative to a reference frame $T_1(O, X_1, Y_1, Z_1)$. Let $T(O, X, Y, Z)$ be the frame attached to the solid such that the $OZ$ axis is the axis of revolution symmetry for the solid.

\begin{enumerate}
    \item What is the number of degrees of freedom of the solid?
    \item Express, in the Resal basis, the rotation vector of $S$ with respect to $T_1$.
    \item Determine the kinetic elements of the solid at point $O$ with respect to $T_1$.
\end{enumerate}

\begin{correctionbox}

\end{correctionbox}

% ----------- EXERCISE 5 -------------
\section{}

A homogeneous and non-deformable square plate $OABC$, of side $a$ and mass $m$, is in motion about point $O$, origin of a reference frame $T_1(O, X_1, Y_1, Z_1)$, such that one of its sides $OA$ remains in contact with the plane $OX_1Y_1$.

\begin{enumerate}
    \item Represent, on a figure, the frame $T(O, X, Y, Z)$ attached to the solid and identify the Euler angles associated with this solid.
    \item What is the number of degrees of freedom of the system?
    \item Determine the inertia operator $J_O$ of the solid at point $O$.
    \item Express, in the basis of $T$, the rotation vector of the solid with respect to $T_1$.
    \item Express, in the basis of $T$, the kinetic torsor of the solid, at $O$, with respect to $T_1$.
    \item Determine the kinetic energy of the solid with respect to $T_1$.
\end{enumerate}

\begin{correctionbox}

\end{correctionbox}

\end{document}